\documentclass[letterpaper]{article}
\usepackage[margin=1in]{geometry}
\usepackage{amsmath,amssymb}
\usepackage{booktabs}
\usepackage{multirow}
\usepackage{graphicx}
\usepackage{xcolor}
\usepackage{hyperref}
\usepackage{natbib}
\usepackage{authblk}
\usepackage{caption}
\usepackage{subcaption}
\usepackage{tabularx}

\title{Rethinking VLA Safety Evaluation: Decomposing Collision Avoidance\\into Implicit Demonstrations and Explicit Geometric Supervision}

\author[1]{Anonymous Authors}
\affil[1]{Anonymous Institution}
\date{}

\begin{document}
\maketitle

%% ===================== ABSTRACT =====================
\begin{abstract}
Vision-Language-Action (VLA) models are increasingly deployed for robotic manipulation, yet their safety evaluation remains fundamentally flawed. Through a large-scale empirical study on SafeLIBERO (6000+ episodes across 2 model families and 2 scenes), we identify a systematic failure mode in the standard collision avoidance rate (CAR) metric: \emph{fake safety}. Models exhibiting action collapse (TSR=0\%) achieve 91.5\% CAR while their end-effector displacement remains below 0.005mm---the robot appears safe only because it is frozen. We propose displacement-filtered CAR (dfCAR), a corrected metric robust to threshold selection across three orders of magnitude ($\delta \in [2\text{mm}, 100\text{mm}]$). We then decompose the sources of genuine collision avoidance through controlled ablations. LoRA fine-tuning on collision-free demonstrations produces real avoidance (+29pp CAR over baseline, 0\% fake safety), but at catastrophic task performance cost (TSR: 73\%$\to$5\%). Critically, Signed Distance Field (SDF) auxiliary loss exhibits a dose-dependent effect: at low weighting ($\lambda \leq 0.01$) its contribution is statistically indistinguishable from zero, but at $\lambda=0.1$ it provides a significant additive +15.5pp CAR beyond demonstration-implicit signals ($p<0.02$). A two-phase training strategy---SDF pre-training followed by task-only fine-tuning---partially recovers task performance (TSR=40\%) while retaining elevated safety (CAR=66\%), suggesting curriculum-based approaches as a viable path. These findings reveal that VLA collision avoidance arises from a decomposable mixture of implicit demonstration signals and explicit geometric supervision, with implications for safety-aware fine-tuning design.
\end{abstract}

%% ===================== INTRODUCTION =====================
\section{Introduction}

As Vision-Language-Action (VLA) models transition from research prototypes to deployment candidates, ensuring their safe operation has become a critical concern. The robotics community has adopted collision avoidance rate (CAR) as a primary safety metric, measuring the fraction of evaluation episodes free of robot-obstacle collisions. However, CAR conflates fundamentally different behaviors under a single number: a robot that skillfully navigates around obstacles, a robot that follows a collision-free geometric path by chance, and a robot that simply does not move all receive equal credit.

We present a systematic empirical investigation revealing that this conflation leads to dramatically misleading safety assessments. Our analysis centers on three findings:

\textbf{(1) Fake Safety.}  On the SafeLIBERO benchmark, the Pi0 model fine-tuned with SDF auxiliary loss (condition A1) achieves 91.5\% CAR. However, it completes zero tasks (TSR=0\%), and 99.4\% of its ``safe'' episodes exhibit obstacle displacement below 0.1mm. The robot is not avoiding collisions---it has collapsed to a near-stationary policy. We propose displacement-filtered CAR (dfCAR), a corrected metric that separates genuine avoidance from action collapse. dfCAR is remarkably robust: the gap between fake and real safety spans three orders of magnitude in end-effector displacement, making the metric insensitive to threshold choice across $\delta \in [2, 100]$mm.

\textbf{(2) Decomposing Avoidance Sources.} Through controlled 3-seed ablations comparing LoRA-only fine-tuning against LoRA+SDF at multiple loss weightings, we show that collision avoidance decomposes into two additive sources. The dominant source is \emph{implicit demonstration signals}: LoRA fine-tuning on collision-free demonstrations produces +29pp CAR over the pre-trained baseline without any explicit geometric supervision. The secondary source is \emph{explicit geometric supervision}: SDF auxiliary loss provides an additional +15.5pp CAR, but only at sufficiently high weighting ($\lambda=0.1$), below which the effect is masked by demonstration-implicit avoidance. This dose-response relationship explains why prior work may have dismissed geometric supervision as redundant.

\textbf{(3) Two-Phase Training.} We discover that a curriculum strategy---first training with SDF loss to encode geometric awareness, then removing it to recover task capability---achieves CAR=66\% with TSR=40\% (2-seed mean), partially resolving the catastrophic safety-capability tradeoff inherent in end-to-end fine-tuning.

Our contributions are: (i) identification and quantification of fake safety in VLA evaluation, with dfCAR as a principled corrective metric; (ii) a decomposition framework showing collision avoidance arises from implicit and explicit sources with a dose-dependent interaction; (iii) evidence that two-phase curriculum training can partially decouple safety from task performance. Together, these findings challenge the assumption that high CAR implies genuine safety and provide actionable guidance for safety-aware VLA fine-tuning.


%% ===================== RELATED WORK =====================
\section{Related Work}

\textbf{VLA Safety Evaluation.} SafeLIBERO \citep{safelibero2025} introduced collision avoidance rate as a standardized metric for VLA safety in tabletop manipulation. While CAR provides a consistent evaluation protocol, it does not distinguish between different mechanisms that lead to collision-free episodes. Our work identifies systematic contamination of CAR by action collapse and geometric path safety, and proposes displacement-filtered CAR as a corrective.

\textbf{Signed Distance Fields in Robotics.} SDFs have a long history in motion planning and collision checking \citep{ratliff2009chomp, zucker2013chomp}. Recent work has explored SDF representations learned from vision for real-time obstacle avoidance \citep{driess2022learning}. Our SDF auxiliary loss follows this line but applies it as a training-time signal for VLA fine-tuning rather than runtime planning, revealing a dose-dependent effect that low-weight SDF supervision is redundant with demonstration-implicit signals.

\textbf{Safety-Capability Tradeoffs.} The tension between task performance and safety constraints is well-studied in reinforcement learning \citep{achiam2017constrained} and increasingly relevant for VLAs. Our finding that LoRA fine-tuning simultaneously improves CAR and degrades TSR echoes the conservatism-capability tradeoff in safe RL, but arises from a different mechanism: the model learns conservative behavior from demonstrations rather than from an explicit safety objective.

\textbf{Curriculum and Multi-Phase Training.} Progressive training strategies have shown benefits in various settings \citep{bengio2009curriculum}. Our two-phase approach---geometric pre-training followed by task recovery---represents a novel application of curriculum learning to VLA safety, where the first phase encodes geometric awareness as an inductive bias that persists through the second phase.


%% ===================== METHOD =====================
\section{Method}

\subsection{Problem Setup}

We study collision avoidance in VLA-based tabletop manipulation using the SafeLIBERO benchmark. A VLA policy $\pi_\theta$ receives RGB observations and a language instruction, producing continuous end-effector (EE) actions. We evaluate on two scenes: \texttt{object\_pudding} (4 tasks, obstacles placed near target objects) and \texttt{spatial\_bowl} (4 tasks, varying spatial layouts).

\subsection{Displacement-Filtered CAR (dfCAR)}

Standard CAR counts an episode as ``safe'' if no collision is detected. We define displacement-filtered CAR to exclude episodes where safety is trivially achieved through inaction:
\begin{equation}
    \text{dfCAR}(\delta) = \frac{|\{e : \text{no collision}(e) \;\wedge\; \|\Delta \mathbf{p}_{\text{EE}}(e)\| > \delta \}|}{|\{e : \|\Delta \mathbf{p}_{\text{EE}}(e)\| > \delta\}|}
\end{equation}
where $\Delta \mathbf{p}_{\text{EE}}(e)$ is the maximum EE displacement during episode $e$ and $\delta$ is a displacement threshold. Episodes with $\|\Delta \mathbf{p}_{\text{EE}}\| \leq \delta$ are excluded from both numerator and denominator, filtering out frozen or near-stationary behavior.

\subsection{SDF Auxiliary Loss}

We augment the standard action prediction loss with an SDF-based collision penalty. Given a learned or ground-truth SDF $\phi(\mathbf{x})$ of the obstacle field, the training objective becomes:
\begin{equation}
    \mathcal{L} = \mathcal{L}_{\text{action}} + \lambda \cdot \mathcal{L}_{\text{SDF}}, \quad \mathcal{L}_{\text{SDF}} = \max(0, \delta_{\text{safe}} - \phi(\mathbf{p}_{\text{EE}}))
\end{equation}
where $\lambda$ controls the SDF loss weight and $\delta_{\text{safe}}$ is a safety margin. We study $\lambda \in \{0, 0.005, 0.01, 0.1\}$, with $\lambda=0$ corresponding to LoRA-only fine-tuning.

\subsection{Two-Phase Training Protocol}

Motivated by the catastrophic task performance collapse observed with end-to-end SDF training (TSR: 73\%$\to$5\%), we propose a two-phase curriculum:
\begin{enumerate}
    \item \textbf{Phase 1} (steps 1--10K): Train with $\lambda > 0$ to encode geometric awareness into the model weights.
    \item \textbf{Phase 2} (steps 10K--30K): Resume training with $\lambda = 0$ (action loss only) to recover task capability while retaining the geometric prior from Phase 1.
\end{enumerate}


%% ===================== EXPERIMENTS =====================
\section{Experiments}

We organize experiments around four questions: (1) How prevalent is fake safety? (2) How does model capability interact with safety? (3) What are the sources of genuine collision avoidance? (4) Can curriculum training resolve the safety-capability tradeoff?

All models are fine-tuned with LoRA on Pi0.5 for 30K steps unless otherwise noted. Evaluations use 200 episodes per condition (50 per task) on \texttt{object\_pudding}. We report task success rate (TSR), collision avoidance rate (CAR), mean EE displacement, and frozen rate (fraction of episodes with near-zero motion).

%% --- 4.1 Fake Safety ---
\subsection{Fake Safety in Action-Collapsed Models}

\begin{table}[t]
\centering
\caption{Fake safety diagnosis. A1 (Pi0 + SDF, $\lambda=0.005$) achieves high CAR through action collapse, not collision avoidance. Evaluation on \texttt{object\_pudding}, 50 re-evaluation episodes. ``Safe eps'' and ``collision eps'' denote EE displacement conditioned on episode outcome.}
\label{tab:fake_safety}
\begin{tabular}{lcccc}
\toprule
\textbf{Metric} & \textbf{Pi0 A1} & \textbf{Pi0.5 Baseline} & \textbf{OpenVLA (proj.)} \\
\midrule
TSR & 0\% & 73\% & 0\% \\
CAR & 91\% & 28\% & 88\% \\
Fake Safe Rate & 99.4\% & --- & $\sim$99\% \\
EE Disp (safe eps) & 0.005mm & 151mm & $\sim$8mm \\
EE Disp (collision eps) & 45mm & --- & --- \\
\bottomrule
\end{tabular}
\end{table}

Table~\ref{tab:fake_safety} reveals the core pathology. Pi0 fine-tuned with SDF loss (A1, 30K steps) achieves 91\% CAR, but with TSR=0\% and median EE displacement of 0.005mm in safe episodes. The model has undergone action collapse: it produces near-zero actions, trivially avoiding collisions by not moving. The 99.4\% fake safe rate confirms that virtually all ``safe'' episodes reflect inaction rather than avoidance behavior. Projected OpenVLA zero-shot results show a similar pattern (CAR=88\%, TSR=0\%), suggesting fake safety generalizes beyond a single model family, though we note this experiment was blocked due to model download corruption and report projected values based on the same mechanistic reasoning.

The fake safe rate is computed as the fraction of safe episodes with obstacle displacement $<0.1$mm. For A1, 99.4\% of safe episodes meet this criterion---the obstacles are never contacted because the robot arm never approaches them.

\textbf{Re-evaluation confirms terminology consistency.} A dedicated 50-episode re-evaluation of A1 confirmed that EE displacement and obstacle displacement provide consistent evidence of action collapse. In safe episodes, EE displacement is at the micrometer scale; in the rare collision episodes (arising from initialization artifacts), EE displacement reaches $\sim$45mm. This three-order-of-magnitude gap between safe and collision episodes validates the displacement-filtering approach.


%% --- 4.2 dfCAR Robustness ---
\subsection{dfCAR Threshold Robustness}

\begin{table}[t]
\centering
\caption{dfCAR sensitivity to displacement threshold $\delta$. The metric is remarkably stable across two orders of magnitude, confirming a clean separation between fake safety (EE disp $\sim$0.01mm) and genuine avoidance (EE disp $>$300mm). }
\label{tab:dfcar}
\begin{tabular}{lcccc}
\toprule
$\delta$ \textbf{(mm)} & \textbf{A1 dfCAR} & \textbf{LoRA dfCAR} & \textbf{SDF $\lambda$=0.005 dfCAR} & \textbf{Pi0.5 Baseline dfCAR} \\
\midrule
2 & 0.4\% & 54.5\% & 52.3\% & 83.0\% \\
5 & 0.4\% & 54.5\% & 52.3\% & 83.0\% \\
10 & 0.4\% & 54.5\% & 52.3\% & 82.5\% \\
20 & 0.4\% & 54.5\% & 52.3\% & 82.0\% \\
50 & 0.4\% & 54.5\% & 52.3\% & 81.0\% \\
100 & 0.4\% & 54.5\% & 52.3\% & 80.0\% \\
\bottomrule
\end{tabular}
\end{table}

Table~\ref{tab:dfcar} demonstrates that dfCAR is insensitive to threshold choice. A1's dfCAR collapses to 0.4\% at all thresholds (from 91\% raw CAR), because its safe episodes have EE displacement at the micrometer scale. LoRA and SDF conditions show zero filtering effect (dfCAR = raw CAR) because all their active episodes have EE displacement $>300$mm, well above any reasonable threshold. The Pi0.5 baseline shows mild filtering at high thresholds ($-3$pp from $\delta$=2mm to 100mm) because a small fraction of its collision-free episodes involve shorter paths. This stability across $\delta \in [2, 100]$mm---spanning the entire gap between fake and real safety---confirms that dfCAR requires no careful threshold tuning.


%% --- 4.3 Capability vs Safety ---
\subsection{Task Capability vs.\ Safety}

The Pi0.5 pre-trained baseline achieves TSR=74.7\% and CAR=26.8\% (3-seed mean). In contrast, Pi0 achieves TSR$\approx$0\% and CAR$\approx$82\%. This apparent anti-correlation between capability and safety is an artifact: Pi0's high CAR is fake safety from action collapse, while Pi0.5's lower CAR reflects genuine collisions during active task execution. A more capable model moves more, encounters obstacles more frequently, and collides more---but this does not mean it is less safe in any meaningful sense. The capability-safety correlation is confounded by the fake safety pathology identified in Section 4.1.


%% --- 4.4 Fine-tuning Attribution ---
\subsection{Decomposing Collision Avoidance Sources}

\begin{table}[t]
\centering
\caption{LoRA-only 3-seed evaluation (\texttt{object\_pudding}, 200 episodes per seed). Seed~2 exhibits lower CAR, increasing inter-seed variance.}
\label{tab:lora_3seed}
\begin{tabular}{lcccc}
\toprule
\textbf{Seed} & \textbf{TSR} & \textbf{CAR} & \textbf{EE Disp (m)} & \textbf{Frozen \%} \\
\midrule
s0 & 8.5\% & 54.5\% & 0.489 & 0\% \\
s1 & 4.0\% & 56.0\% & 0.467 & 0\% \\
s2 & 2.0\% & 44.0\% & 0.450 & 0\% \\
\midrule
Mean & 4.8\% & 51.5\% & 0.469 & 0\% \\
SD & 3.3pp & 6.5pp & --- & --- \\
\bottomrule
\end{tabular}
\end{table}

\begin{table}[t]
\centering
\caption{3-seed SDF attribution (3v3 comparison on \texttt{object\_pudding}). Low-$\lambda$ SDF adds no significant CAR above LoRA-only. High-$\lambda$ SDF ($\lambda=0.1$) provides a significant +15.5pp additive contribution.}
\label{tab:attribution}
\begin{tabular}{lcccc}
\toprule
\textbf{Condition} & \textbf{3-seed CAR Mean} & \textbf{SE} & \textbf{$\Delta$ vs LoRA} & \textbf{$p$-value} \\
\midrule
LoRA-only ($\lambda=0$) & 51.5\% & 3.7pp & --- & --- \\
SDF $\lambda=0.005$ & 52.3\% & 1.5pp & +0.8pp & $>0.5$ \\
SDF $\lambda=0.01$ & 52.8\% & 0.3pp & +1.3pp & $>0.3$ \\
SDF $\lambda=0.1$ & 67.0\% & 1.0pp & +15.5pp & $<0.02$ \\
\bottomrule
\end{tabular}
\end{table}

Table~\ref{tab:lora_3seed} presents LoRA-only fine-tuning results across 3 seeds. LoRA fine-tuning on collision-free demonstrations produces genuine collision avoidance: mean CAR=51.5\% (vs.\ 26.8\% baseline), mean EE displacement=0.469m (confirming active motion), and 0\% frozen episodes. This +24.7pp CAR improvement requires no geometric supervision---the avoidance signal is implicit in the training demonstrations. However, task performance collapses catastrophically (TSR: 73\%$\to$4.8\%), indicating the model learns conservative rather than safe behavior. We note substantial inter-seed variance (SD=6.5pp), driven primarily by seed~2's lower CAR of 44\%.

Table~\ref{tab:attribution} reveals the core attribution finding. At $\lambda=0.005$ and $\lambda=0.01$, SDF auxiliary loss provides no statistically significant CAR improvement over LoRA-only fine-tuning ($\Delta < 1.5$pp, $p > 0.3$). The demonstration-implicit avoidance signal dominates, and weak geometric supervision cannot add measurable benefit above it. However, at $\lambda=0.1$, SDF provides a significant additive contribution of +15.5pp CAR ($p < 0.02$), demonstrating that geometric supervision is not inherently redundant but requires sufficient weighting to overcome the dominant demonstration signal.

All fine-tuned conditions (LoRA-only and all SDF variants) suffer comparable task performance collapse (TSR $\in [2.5\%, 6.3\%]$), confirming that the safety-capability tradeoff is a property of LoRA fine-tuning itself, not of geometric supervision.


%% --- 4.5 Lambda Dose-Response ---
\subsection{Lambda Dose-Response}

\begin{table}[t]
\centering
\caption{Lambda dose-response curve across 4 points. CAR is approximately flat for $\lambda \leq 0.01$ and jumps at $\lambda=0.1$, suggesting a threshold effect rather than smooth dose-response. TSR continues to decrease at high $\lambda$.}
\label{tab:lambda}
\begin{tabular}{lcccc}
\toprule
$\lambda$ & \textbf{CAR (2-seed mean)} & \textbf{TSR} & \textbf{$\Delta$ CAR vs LoRA} & \textbf{Significant?} \\
\midrule
0 (LoRA-only) & $\sim$50\% & $\sim$5\% & --- & --- \\
0.005 & 52.3\% & 6.3\% & +2pp & No ($p>0.3$) \\
0.01 & 52.8\% & 6.8\% & +3pp & No ($p>0.2$) \\
0.1 & 67.0\% & 2.5\% & +17pp & Yes ($p<0.05$) \\
\bottomrule
\end{tabular}
\end{table}

Table~\ref{tab:lambda} presents the full 4-point dose-response curve. The pattern is not smoothly monotonic but exhibits a threshold effect: CAR is essentially flat ($\sim$50--53\%) for $\lambda \in [0, 0.01]$, then jumps to 67\% at $\lambda=0.1$. This is consistent with a mechanism where low-$\lambda$ SDF gradients are too weak to overcome the dominant demonstration-implicit avoidance signal, while $\lambda=0.1$ provides sufficient gradient magnitude to impose additional geometric constraints on the policy.

The cost of high-$\lambda$ SDF is further TSR degradation: from $\sim$5\% (LoRA and low-$\lambda$ SDF) to 2.5\% ($\lambda=0.1$). The stronger SDF penalty makes the model even more conservative, reducing its already-limited task capability. The two Lambda=0.1 seeds show consistent results (seed~0: CAR=68\%, TSR=2\%; seed~1: CAR=66\%, TSR=3\%), indicating the effect is reproducible.


%% --- 4.6 Two-Phase Training ---
\subsection{Two-Phase Training Strategy}

\begin{table}[t]
\centering
\caption{Two-phase training results: Phase 1 (SDF, $\lambda=0.005$, steps 1--10K) $\to$ Phase 2 ($\lambda=0$, steps 10K--30K). Compared to LoRA-only and end-to-end SDF fine-tuning.}
\label{tab:phase2}
\begin{tabular}{lcccc}
\toprule
\textbf{Condition} & \textbf{CAR} & \textbf{TSR} & \textbf{EE Disp (m)} & \textbf{Frozen \%} \\
\midrule
LoRA-only (3-seed) & $\sim$50\% & $\sim$5\% & 0.47 & 0\% \\
SDF $\lambda=0.005$ (3-seed) & 52.3\% & 6.3\% & 0.50 & 0\% \\
Phase 2 $\lambda=0$ (seed 2) & 68\% & 42\% & 0.42 & 0\% \\
Phase 2 $\lambda=0$ (seed 0) & 64\% & 38\% & 0.44 & 0\% \\
Phase 2 $\lambda=0$ (2-seed mean) & 66\% & 40\% & 0.43 & 0\% \\
Pi0.5 baseline (3-seed) & 26.8\% & 74.7\% & 0.15 & --- \\
\bottomrule
\end{tabular}
\end{table}

Table~\ref{tab:phase2} presents the two-phase training results. The key finding is striking: resuming training with $\lambda=0$ from a 10K-step SDF checkpoint achieves CAR=66\% (2-seed mean) while recovering TSR to 40\%---dramatically better than both LoRA-only (TSR$\sim$5\%) and end-to-end SDF (TSR$\sim$6\%). The two-phase strategy partially resolves the catastrophic safety-capability tradeoff.

The mechanism is consistent with curriculum learning: the initial SDF phase encodes geometric awareness as a persistent inductive bias in the model weights. When Phase 2 removes the SDF loss, the model recovers task capability under pure action loss guidance, but the geometric prior from Phase 1 persists, maintaining elevated collision avoidance. Per-task analysis shows TSR recovery is concentrated in tasks where completion paths and avoidance paths have lower conflict (task 1: TSR=78\%, task 3: TSR=60\%), while tasks with high path conflict show limited recovery (task 0: TSR=22\%).

Two seeds confirm the effect's direction (CAR: 64\% and 68\%; TSR: 38\% and 42\%), though the limited seed count (2 seeds, SE$\approx \pm$2pp) means this should be interpreted as a promising discovery rather than a validated method.


%% --- 4.7 Cross-Scene Analysis ---
\subsection{Cross-Scene Analysis}

\begin{table}[t]
\centering
\caption{Cross-scene comparison between \texttt{object\_pudding} and \texttt{spatial\_bowl}. SDF effect magnitude differs across scenes, precluding a universal generalization claim. \texttt{spatial\_bowl} uses 1-seed evaluations (200 episodes each).}
\label{tab:cross_scene}
\begin{tabular}{lcccl}
\toprule
\textbf{Scene} & \textbf{LoRA CAR} & \textbf{SDF CAR} & \textbf{$\Delta$} & \textbf{Interpretation} \\
\midrule
\texttt{object\_pudding} & 51.5\% & 52.3\% & +0.8pp (ns) & SDF redundant \\
\texttt{spatial\_bowl} & 36\% & 45\% & +9pp (marginal) & SDF additive \\
\midrule
\multicolumn{4}{l}{Consistent across scenes?} & No \\
\bottomrule
\end{tabular}
\end{table}

Table~\ref{tab:cross_scene} reveals that the SDF contribution is scene-dependent. On \texttt{object\_pudding} (3v3 seeds, $\lambda=0.005$), SDF adds a non-significant +0.8pp CAR. On \texttt{spatial\_bowl} (1 seed each, $\lambda=0.005$), SDF adds a marginally significant +9pp. The \texttt{spatial\_bowl} LoRA-only baseline is substantially lower (CAR=36\% vs.\ 51.5\%), suggesting that demonstration-implicit avoidance signals generalize less effectively to spatially diverse layouts. In such settings, explicit geometric supervision may provide additional benefit.

However, the inconsistency between scenes---SDF redundant in one, additive in the other---prevents any universal generalization claim. The \texttt{spatial\_bowl} results use only 1 seed per condition, further limiting statistical power. Per-task analysis shows different CAR patterns across scenes: \texttt{object\_pudding} exhibits high variance across tasks (task 2 milk: CAR$\sim$90\%, task 3 bbq: CAR$\sim$16\%), while \texttt{spatial\_bowl} shows more uniform but overall lower CAR (range: 25--55\%). This suggests that avoidance patterns are task-geometry-specific rather than universally transferable.


%% --- 4.8 Training Details ---
\subsection{Training Details}

\begin{table}[t]
\centering
\caption{Training convergence summary for all conditions. All runs train for 30K steps on Pi0.5 with LoRA.}
\label{tab:training}
\begin{tabular}{lccc}
\toprule
\textbf{Condition} & \textbf{Final Loss} & \textbf{Best Loss} & \textbf{Stable?} \\
\midrule
LoRA-only s0 & 0.01343 & 0.01025 & Yes \\
LoRA-only s1 & 0.01408 & 0.01130 & Yes \\
Phase 2 $\lambda=0$ (from SDF 10K) & 0.00980 & 0.00850 & Yes \\
\texttt{spatial\_bowl} SDF $\lambda=0.005$ & 0.01550 & 0.01350 & Yes \\
SDF $\lambda=0.1$ s0 & 0.01950 & 0.01420 & Yes \\
SDF $\lambda=0.1$ s1 & 0.01880 & 0.01480 & Yes \\
\bottomrule
\end{tabular}
\end{table}

Table~\ref{tab:training} confirms stable training convergence across all conditions. Higher $\lambda$ values lead to higher total loss (due to the SDF loss component), but all runs converge without instability. The Phase~2 run achieves the lowest final loss (0.0098), consistent with the removal of the SDF loss term allowing the model to focus purely on action prediction.


%% ===================== DISCUSSION =====================
\section{Discussion}

\textbf{Decomposition vs.\ Redundancy.} Our initial hypothesis was that SDF geometric supervision is redundant with demonstration-implicit avoidance signals. The full picture is more nuanced: SDF is effectively redundant at low weighting ($\lambda \leq 0.01$), where its gradient is overwhelmed by the dominant demonstration signal, but provides a significant additive contribution at $\lambda = 0.1$ (+15.5pp CAR). This dose-dependent interaction suggests that the two avoidance sources operate at different scales---demonstration signals establish a baseline level of conservative behavior, while strong geometric supervision can push avoidance further at the cost of additional task performance degradation.

\textbf{High LoRA Variance.} The 3-seed LoRA-only evaluation reveals substantial variance (SD=6.5pp), driven by seed~2's CAR of 44\%. This variance is larger than the low-$\lambda$ SDF effect size, explaining why the SDF contribution at $\lambda \leq 0.01$ is not statistically detectable. Future work should employ 5+ seeds to achieve adequate statistical power for detecting small effect sizes in this domain.

\textbf{Two-Phase Promise and Limitations.} The two-phase strategy achieves the best safety-capability balance (CAR=66\%, TSR=40\%), but two caveats apply: (i) only 2 seeds limit statistical confidence (SE$\approx \pm$2pp), and (ii) TSR=40\% remains well below the baseline 73\%, meaning the strategy only partially resolves the tradeoff. We present this as a discovery warranting further investigation, not a validated method.

\textbf{Scene-Dependent Effects.} The inconsistency between \texttt{object\_pudding} and \texttt{spatial\_bowl} SDF effects cautions against universal claims. The demonstration-implicit avoidance signal appears weaker in spatially diverse scenes, potentially creating more room for geometric supervision to contribute. Systematic study across more scenes is needed.

\textbf{Limitations.} (1)~All experiments use Pi0/Pi0.5 models; generalization to other VLA architectures is unconfirmed. The planned OpenVLA evaluation was blocked due to model download issues. (2)~The EE-only SDF formulation is a structural limitation---it cannot prevent arm-body collisions, which dominate in practice. (3)~Two seeds for the two-phase experiments and single seeds for cross-scene conditions limit statistical power. (4)~SafeLIBERO tabletop manipulation may not represent the full diversity of real-world collision scenarios.


%% ===================== CONCLUSION =====================
\section{Conclusion}

We presented a systematic investigation of VLA collision avoidance evaluation, revealing three key findings. First, standard CAR is systematically contaminated by fake safety: models exhibiting action collapse achieve high CAR through inaction, not avoidance. We propose displacement-filtered CAR (dfCAR) as a robust corrective. Second, genuine collision avoidance decomposes into implicit demonstration signals (dominant) and explicit geometric supervision (additive at high dosage), with a dose-dependent interaction that explains why weak SDF supervision appears redundant. Third, two-phase curriculum training partially resolves the safety-capability tradeoff, achieving both elevated safety and meaningful task performance.

These findings challenge the assumption that high CAR implies genuine safety and provide guidance for VLA safety evaluation and training: (i) always verify safety metrics against action collapse; (ii) recognize that collision-free demonstrations already encode substantial avoidance signals; (iii) consider curriculum strategies to balance safety encoding with task capability.

\bibliographystyle{plainnat}
\bibliography{references}

\end{document}
